% arara: xelatex: { options: ["--synctex=1", "-interaction=nonstopmode"] }
% arara: biber
% arara: xelatex: { options: ['-synctex=1', '-interaction=nonstopmode'] }
% arara: xelatex: { options: ['-synctex=1', '-interaction=nonstopmode'] }

\documentclass[aspectratio=169,10pt]{beamer}
\usepackage{graphicx}
\usepackage{amssymb,amsfonts,amsmath}
\usepackage{array}
\usepackage{xeCJK}
\usepackage{tcolorbox}
\usepackage{tikz}
\usetikzlibrary{positioning,shapes.geometric,arrows.meta,calc}
\usepackage{pifont}
\setCJKmainfont{STSong}
\setCJKsansfont{STHeiti}
\setCJKmonofont{STFangsong}
\usefonttheme[onlymath]{serif}

% 自定义颜色
\definecolor{highlightyellow}{RGB}{255, 242, 0}
\definecolor{dotred}{RGB}{200, 30, 30}
\definecolor{dotyellow}{RGB}{255, 200, 0}
\definecolor{darkgreen}{RGB}{0, 128, 0}
\definecolor{accentblue}{RGB}{21,101,192}

% biblatex with biber
\usepackage[backend=biber,style=numeric,sorting=none,firstinits=true]{biblatex}
\addbibresource{ref.bib}
\usepackage{hyperref}

\title{开发 Coding Agent 的技术要点\\{\large 基于 Context Graph 的研究洞察}}
\author{WU Zihan}
\date{\today}

\begin{document}
\maketitle

\begin{frame}
  \frametitle{目录}
  \tableofcontents
\end{frame}

%==============================================================================
\section{核心洞察总结}
%==============================================================================

\begin{frame}
  \frametitle{核心洞察 1：Agent 的根本问题是"记忆"}

  \begin{columns}[T]
    \begin{column}{0.55\textwidth}
      \begin{block}{技术演进路径}
        \small
        LLM $\xrightarrow{+工具}$ Agent $\xrightarrow{+记忆}$ Agent+Memory $\xrightarrow{+图结构}$ Agent+Context Graph
      \end{block}

      \vspace{0.3cm}

      \begin{alertblock}{四大核心问题}
        \footnotesize
        \begin{enumerate}
          \item \textbf{重复犯错}：同样的错误反复出现
          \item \textbf{无法积累经验}：每次从零开始
          \item \textbf{上下文爆炸}：历史记录过长导致性能下降
          \item \textbf{跨任务学习困难}：经验无法迁移
        \end{enumerate}
      \end{alertblock}
    \end{column}

    \begin{column}{0.42\textwidth}
      \begin{block}{解决方案}
        \textbf{Context Graph}：结构化记忆系统

        \vspace{0.2cm}

        \footnotesize
        \begin{itemize}
          \item \textcolor{darkgreen}{优于 RAG}：保留实体间关系，支持多跳推理
          \item \textcolor{darkgreen}{优于传统 KG}：包含时间、来源等上下文信息
          \item \textcolor{darkgreen}{可学习}：根据反馈动态优化
        \end{itemize}
      \end{block}
    \end{column}
  \end{columns}
\end{frame}

\begin{frame}
  \frametitle{核心洞察 2：两条技术路线殊途同归}

  \begin{columns}[T]
    \begin{column}{0.48\textwidth}
      \begin{block}{路线 1：图增强检索（GraphRAG）}
        \footnotesize
        \textbf{目标}：改进 LLM 检索质量

        \vspace{0.2cm}

        \textbf{代表工作}：
        \begin{itemize}
          \item GraphRAG (MS)：社区检测+分层摘要
          \item HippoRAG：海马体索引理论
          \item CGR³：上下文增强知识图谱
        \end{itemize}

        \vspace{0.2cm}

        \textcolor{darkgreen}{成熟度高}：2024-2025 学术验证充分
      \end{block}
    \end{column}

    \begin{column}{0.48\textwidth}
      \begin{block}{路线 2：Agent Memory}
        \footnotesize
        \textbf{目标}：构建可学习的经验系统

        \vspace{0.2cm}

        \textbf{代表工作}：
        \begin{itemize}
          \item ExpeL：提取洞察+检索经验
          \item A-Mem：动态自组织记忆网络
          \item Trainable Graph：多层图+RL优化
        \end{itemize}

        \vspace{0.2cm}

        \textcolor{red}{Coding 场景几乎空白}
      \end{block}
    \end{column}
  \end{columns}

  \vspace{0.3cm}

  \begin{alertblock}{关键发现}
    \textbf{GraphRAG 是检索技术，Context Graph 是架构愿景}——两者是同一技术栈的不同层次
  \end{alertblock}
\end{frame}

\begin{frame}
  \frametitle{核心洞察 3：图结构的五大优势}

  \small
  \renewcommand{\arraystretch}{1.4}
  \centering
  \begin{tabular}{>{\centering}p{2.2cm}|>{\centering}p{4.5cm}|>{\centering\arraybackslash}p{4.8cm}}
    \textbf{能力} & \textbf{RAG 的局限} & \textbf{Graph 的优势} \\
    \hline
    关系建模 & 找相似文档，可能不相关 & 出边直接指向下一步，精确 \\
    \hline
    路径推理 & 无"路径"概念 & 搜索最短/最优路径 \\
    \hline
    失败诊断 & 找类似案例，不知差异 & 比较成功/失败路径，找分叉点 \\
    \hline
    动态更新 & Fine-tuning 需重训，可能遗忘 & 增量更新节点/边 \\
    \hline
    可解释性 & 神经网络黑盒 & "因为 A→B 成功率 95\%" \\
  \end{tabular}
\end{frame}

\begin{frame}
  \frametitle{核心洞察 4：Coding Agent 是理想应用场景}

  \begin{columns}[T]
    \begin{column}{0.58\textwidth}
      \begin{block}{为什么 Coding 是理想场景？}
        \footnotesize
        \begin{enumerate}
          \item \textbf{反馈明确}：测试通过/失败、编译成功/失败
          \item \textbf{结构化强}：代码有明确的依赖关系
          \item \textbf{可重复}：相同输入可多次测试
          \item \textbf{经验可迁移}：Bug 模式、修复策略有共性
          \item \textbf{成本高昂}：token 消耗大，需要优化
        \end{enumerate}
      \end{block}

      \begin{alertblock}{研究空白}
        \small
        现有工作主要在通用 QA 任务验证，\\
        \textbf{Coding 场景几乎未探索}
      \end{alertblock}
    \end{column}

    \begin{column}{0.38\textwidth}
      \begin{block}{主要任务类型}
        \footnotesize
        \begin{itemize}
          \item 代码生成
          \item \textbf{Bug 修复 (SWE-bench)}
          \item 漏洞检测
          \item 测试生成
          \item 代码审查
          \item 代码理解
        \end{itemize}
      \end{block}

      \vspace{0.2cm}

      \begin{block}{推荐起点}
        \footnotesize
        \textbf{Bug 修复}：
        \begin{itemize}
          \item 成熟的 benchmark
          \item 明确的评估标准
          \item 与漏洞检测相关
        \end{itemize}
      \end{block}
    \end{column}
  \end{columns}
\end{frame}

%==============================================================================
\section{开发 Coding Agent 的技术要点}
%==============================================================================

\begin{frame}
  \frametitle{技术架构全景}

  \begin{center}
  \begin{tikzpicture}[scale=0.9, every node/.style={font=\small}]
    % 定义颜色
    \definecolor{layer1}{RGB}{227,242,253}
    \definecolor{layer2}{RGB}{187,222,251}
    \definecolor{layer3}{RGB}{144,202,249}
    \definecolor{layer4}{RGB}{66,165,245}

    % 基础层
    \draw[fill=layer1, rounded corners] (-6, -1) rectangle (6, 0);
    \node at (0, -0.5) {\textbf{基础层}：LLM (GPT-4, Claude, DeepSeek)};

    % 工具层
    \draw[fill=layer2, rounded corners] (-6, 0.2) rectangle (6, 1.2);
    \node at (0, 0.7) {\textbf{工具层}：代码执行、测试运行、文件操作、搜索};

    % 记忆层
    \draw[fill=layer3, rounded corners] (-6, 1.4) rectangle (6, 2.4);
    \node at (0, 1.9) {\textbf{记忆层}：Context Graph（本报告重点）};

    % 任务层
    \draw[fill=layer4, rounded corners] (-6, 2.6) rectangle (6, 3.6);
    \node[text=white] at (0, 3.1) {\textbf{任务层}：Bug 修复、代码生成、测试、审查};

    % 箭头
    \draw[->, thick] (0, 0) -- (0, 0.2);
    \draw[->, thick] (0, 1.2) -- (0, 1.4);
    \draw[->, thick] (0, 2.4) -- (0, 2.6);
  \end{tikzpicture}
  \end{center}

  \vspace{0.3cm}

  \begin{alertblock}{关键点}
    \textbf{记忆层是差异化的核心}——基础层和工具层都是通用技术，记忆层决定了 Agent 的学习能力
  \end{alertblock}
\end{frame}

\begin{frame}
  \frametitle{工作 1：设计 Context Graph 数据结构}

  \begin{columns}[T]
    \begin{column}{0.48\textwidth}
      \begin{block}{节点类型设计}
        \footnotesize
        \textbf{1. Action Node（动作节点）}
        \begin{itemize}
          \item \texttt{action\_type}: read\_file, write\_file, run\_test
          \item \texttt{parameters}: 参数
          \item \texttt{result}: success/failure
          \item \texttt{cost}: token 消耗
        \end{itemize}

        \vspace{0.2cm}

        \textbf{2. State Node（状态节点）}
        \begin{itemize}
          \item \texttt{description}: "发现 bug 在 line 42"
          \item \texttt{context}: 相关文件、变量
          \item \texttt{is\_terminal}: 是否终止
        \end{itemize}

        \vspace{0.2cm}

        \textbf{3. Insight Node（洞察节点）}
        \begin{itemize}
          \item \texttt{insight}: 提取的经验规则
          \item \texttt{confidence}: 置信度
          \item \texttt{source\_tasks}: 来源任务
        \end{itemize}
      \end{block}
    \end{column}

    \begin{column}{0.48\textwidth}
      \begin{block}{边类型设计}
        \footnotesize
        \textbf{1. NEXT（时序/因果）}
        \begin{itemize}
          \item Action\_A $\rightarrow$ Action\_B
          \item 属性: success\_rate, count
        \end{itemize}

        \vspace{0.2cm}

        \textbf{2. CAUSES（因果）}
        \begin{itemize}
          \item Action\_X $\rightarrow$ Failure\_Y
          \item 失败原因追踪
        \end{itemize}

        \vspace{0.2cm}

        \textbf{3. SIMILAR\_TO（相似）}
        \begin{itemize}
          \item Task\_1 $\leftrightarrow$ Task\_2
          \item 相似任务经验检索
        \end{itemize}

        \vspace{0.2cm}

        \textbf{4. ABSTRACTS（抽象）}
        \begin{itemize}
          \item [Actions] $\rightarrow$ Insight
          \item 具体经验 $\rightarrow$ 通用规则
        \end{itemize}
      \end{block}
    \end{column}
  \end{columns}
\end{frame}

\begin{frame}
  \frametitle{工作 2：实现图的构建与更新机制}

  \begin{block}{图构建流程}
    \small
    \begin{enumerate}
      \item \textbf{轨迹记录}：Agent 每次执行任务时，记录所有 action 和状态转换
      \item \textbf{节点创建}：将 action/state 转化为图节点，提取关键属性
      \item \textbf{边连接}：建立节点间的关系（时序、因果、相似）
      \item \textbf{洞察提取}：用 LLM 从成功/失败轨迹中提取高层规则
      \item \textbf{图整合}：将新子图合并到全局图中
    \end{enumerate}
  \end{block}

  \vspace{0.2cm}

  \begin{columns}[T]
    \begin{column}{0.48\textwidth}
      \begin{block}{技术挑战 1：动态更新}
        \footnotesize
        \begin{itemize}
          \item 如何增量更新而非重建？
          \item 如何处理冲突信息？
          \item 如何遗忘过时信息？
        \end{itemize}
      \end{block}
    \end{column}

    \begin{column}{0.48\textwidth}
      \begin{block}{技术挑战 2：存储效率}
        \footnotesize
        \begin{itemize}
          \item 图数据库选择（Neo4j, NetworkX）
          \item 大规模图的索引优化
          \item 分布式存储策略
        \end{itemize}
      \end{block}
    \end{column}
  \end{columns}
\end{frame}

\begin{frame}
  \frametitle{工作 3：实现检索与推理算法}

  \begin{columns}[T]
    \begin{column}{0.55\textwidth}
      \begin{block}{检索策略}
        \footnotesize
        \textbf{1. 相似任务检索}
        \begin{itemize}
          \item 输入：当前 Bug 描述
          \item 算法：embedding 相似度 + 图结构相似度
          \item 输出：Top-K 相似任务的成功路径
        \end{itemize}

        \vspace{0.2cm}

        \textbf{2. 相关洞察检索}
        \begin{itemize}
          \item 输入：当前状态、已执行动作
          \item 算法：基于关键词 + 图遍历
          \item 输出：适用的经验规则
        \end{itemize}

        \vspace{0.2cm}

        \textbf{3. 失败路径过滤}
        \begin{itemize}
          \item 排除已知失败的 action 序列
          \item 避免重复犯错
        \end{itemize}
      \end{block}
    \end{column}

    \begin{column}{0.42\textwidth}
      \begin{block}{推理算法}
        \footnotesize
        \textbf{路径规划}
        \begin{itemize}
          \item 从当前状态到目标状态
          \item 考虑成功率和 cost
          \item A* / Dijkstra 变体
        \end{itemize}

        \vspace{0.2cm}

        \textbf{失败诊断}
        \begin{itemize}
          \item 比较成功/失败路径
          \item 找到关键分叉点
          \item 定位失败原因
        \end{itemize}
      \end{block}

      \vspace{0.2cm}

      \begin{alertblock}{核心}
        \tiny
        \textbf{图+LLM 协同}：图提供结构化检索，LLM 负责语义理解
      \end{alertblock}
    \end{column}
  \end{columns}
\end{frame}

\begin{frame}[fragile]
  \frametitle{工作 4：构建任务执行主循环}

  \begin{columns}[T]
    \begin{column}{0.58\textwidth}
      \begin{block}{主循环伪代码}
        \footnotesize
        \begin{verbatim}
while not task_completed:
  # 1. 检索相关经验
  similar_tasks = graph.retrieve_similar(task)
  insights = graph.retrieve_insights(current_state)

  # 2. 规划下一步动作
  action = planner.decide(
      current_state,
      similar_tasks,
      insights
  )

  # 3. 执行动作
  result = executor.run(action)

  # 4. 更新图
  graph.add_node(action, result)
  graph.add_edge(prev_action, action)

  # 5. 检查终止条件
  if result.is_terminal:
    break
        \end{verbatim}
      \end{block}
    \end{column}

    \begin{column}{0.38\textwidth}
      \begin{block}{关键组件}
        \footnotesize
        \begin{enumerate}
          \item \textbf{Planner}：决策模块
          \begin{itemize}
            \item 输入：状态+经验
            \item 输出：下一步 action
          \end{itemize}

          \vspace{0.2cm}

          \item \textbf{Executor}：执行模块
          \begin{itemize}
            \item 调用工具（测试、编译）
            \item 解析结果
          \end{itemize}

          \vspace{0.2cm}

          \item \textbf{Graph Manager}
          \begin{itemize}
            \item 图的增删改查
            \item 洞察提取
          \end{itemize}
        \end{enumerate}
      \end{block}
    \end{column}
  \end{columns}
\end{frame}

\begin{frame}
  \frametitle{工作 5：实现学习与优化机制}

  \begin{block}{学习目标}
    让 Agent 从每次任务中学习，持续提升成功率和效率
  \end{block}

  \begin{columns}[T]
    \begin{column}{0.48\textwidth}
      \begin{block}{方法 1：基于规则的学习}
        \footnotesize
        \textbf{洞察提取}（ExpeL 方法）
        \begin{enumerate}
          \item 收集成功和失败轨迹
          \item 用 LLM 提取共性规则
          \item 存储为 Insight Node
        \end{enumerate}

        \vspace{0.2cm}

        \textbf{优点}：可解释性强\\
        \textbf{缺点}：依赖 LLM 质量
      \end{block}

      \vspace{0.2cm}

      \begin{block}{方法 2：基于统计的学习}
        \footnotesize
        \textbf{边权重更新}
        \begin{itemize}
          \item 记录每条边的成功次数/总次数
          \item 计算 success\_rate
          \item 规划时优先选择高成功率路径
        \end{itemize}
      \end{block}
    \end{column}

    \begin{column}{0.48\textwidth}
      \begin{block}{方法 3：强化学习（进阶）}
        \footnotesize
        \textbf{可训练图记忆}方法
        \begin{enumerate}
          \item 定义奖励：任务成功 +1，失败 -1
          \item 训练策略网络：state $\rightarrow$ action
          \item 更新图中的策略权重
        \end{enumerate}

        \vspace{0.2cm}

        \textbf{挑战}：
        \begin{itemize}
          \item 需要大量训练数据
          \item 如何定义状态空间？
          \item 探索-利用平衡
        \end{itemize}
      \end{block}
    \end{column}
  \end{columns}

  \vspace{0.2cm}

  \begin{alertblock}{推荐}
    \small
    初期：方法 1 + 方法 2 结合；长期：探索方法 3
  \end{alertblock}
\end{frame}

\begin{frame}
  \frametitle{工作 6：设计评估与基准测试}

  \begin{columns}[T]
    \begin{column}{0.55\textwidth}
      \begin{block}{评估指标体系}
        \footnotesize
        \textbf{1. 任务成功率}
        \begin{itemize}
          \item Resolved Rate (SWE-bench)
          \item Pass@K (代码生成)
        \end{itemize}

        \vspace{0.2cm}

        \textbf{2. 效率指标}
        \begin{itemize}
          \item Token 消耗
          \item 执行步数
          \item 运行时间
        \end{itemize}

        \vspace{0.2cm}

        \textbf{3. 学习能力}
        \begin{itemize}
          \item 首次尝试 vs 第 N 次尝试的成功率
          \item 跨任务迁移效果
        \end{itemize}

        \vspace{0.2cm}

        \textbf{4. 图质量}
        \begin{itemize}
          \item 图的规模（节点数、边数）
          \item 检索精度/召回率
        \end{itemize}
      \end{block}
    \end{column}

    \begin{column}{0.42\textwidth}
      \begin{block}{基准数据集}
        \footnotesize
        \textbf{Bug 修复}
        \begin{itemize}
          \item SWE-bench
          \item SWE-bench Lite
        \end{itemize}

        \vspace{0.2cm}

        \textbf{代码生成}
        \begin{itemize}
          \item HumanEval
          \item MBPP
        \end{itemize}

        \vspace{0.2cm}

        \textbf{漏洞检测}
        \begin{itemize}
          \item CodeMender
          \item Spark
        \end{itemize}
      \end{block}

      \vspace{0.2cm}

      \begin{alertblock}{对比基线}
        \tiny
        \begin{itemize}
          \item 无记忆 Agent
          \item RAG Agent
          \item ExpeL / A-Mem
        \end{itemize}
      \end{alertblock}
    \end{column}
  \end{columns}
\end{frame}

\begin{frame}
  \frametitle{工作 7：处理 Coding 场景的特殊挑战}

  \begin{block}{挑战 1：代码库动态变化}
    \small
    \textbf{问题}：代码更新后，历史经验可能失效\\
    \textbf{解决}：
    \begin{itemize}
      \item 为每个经验记录代码版本/commit hash
      \item 检索时检查代码是否仍然有效
      \item 定期清理过时经验
    \end{itemize}
  \end{block}

  \begin{columns}[T]
    \begin{column}{0.48\textwidth}
      \begin{block}{挑战 2：Bug 相似度判断}
        \footnotesize
        \textbf{问题}：如何判断两个 Bug 是否相似？

        \vspace{0.1cm}

        \textbf{特征}：
        \begin{itemize}
          \item 错误类型（语法/逻辑/性能）
          \item 涉及的代码模块
          \item 错误堆栈相似度
          \item Issue 描述 embedding
        \end{itemize}

        \vspace{0.1cm}

        \textbf{算法}：混合相似度计算
      \end{block}
    \end{column}

    \begin{column}{0.48\textwidth}
      \begin{block}{挑战 3：代码结构理解}
        \footnotesize
        \textbf{问题}：需要理解项目依赖关系

        \vspace{0.1cm}

        \textbf{方法}：
        \begin{itemize}
          \item 构建代码依赖图（AST 分析）
          \item 与 Context Graph 融合
          \item 检索时考虑代码结构
        \end{itemize}

        \vspace{0.1cm}

        \textbf{参考}：Think-on-Graph 的交互式探索
      \end{block}
    \end{column}
  \end{columns}
\end{frame}

%==============================================================================
\section{业界实践与趋势}
%==============================================================================

\begin{frame}
  \frametitle{主流 AI 产品的记忆机制对比}

  \tiny
  \renewcommand{\arraystretch}{1.3}
  \begin{tabular}{>{\centering}p{1.5cm}|>{\centering}p{2.2cm}|>{\centering}p{2.2cm}|>{\centering}p{2.2cm}|>{\centering\arraybackslash}p{2cm}}
    \textbf{产品} & \textbf{短期记忆} & \textbf{长期记忆} & \textbf{存储机制} & \textbf{图谱化} \\
    \hline
    ChatGPT & 32K tokens 窗口 & 全历史检索 (2025) & 向量索引 + 用户档案 & 否 \\
    \hline
    Claude & 200K tokens 窗口 & 项目级知识库 (Markdown) & RAG + 文件系统 & 部分（关系推理） \\
    \hline
    Gemini & 2M tokens 窗口 & Personal Intelligence 跨应用 & 动态上下文打包 + 多源检索 & 是（个人知识图谱） \\
    \hline
    Pi & 数千 tokens & 持久用户画像 + 情感记忆 & 云端数据库 + embedding & 否（但有关系上下文） \\
  \end{tabular}

  \vspace{0.3cm}

  \begin{block}{关键趋势}
    \footnotesize
    \begin{itemize}
      \item \textbf{上下文窗口竞赛}：从 4K → 200K → 2M tokens
      \item \textbf{跨会话记忆}：从隔离对话到持续学习
      \item \textbf{个性化深化}：从通用助手到私人伴侣
      \item \textbf{图谱化萌芽}：Gemini 率先整合个人知识图谱
    \end{itemize}
  \end{block}
\end{frame}

\begin{frame}
  \frametitle{开源框架的记忆实践}

  \begin{columns}[T]
    \begin{column}{0.48\textwidth}
      \begin{block}{MemGPT/Letta}
        \footnotesize
        \textbf{核心理念}："LLM 作为操作系统"

        \vspace{0.2cm}

        \textbf{创新点}：
        \begin{itemize}
          \item 虚拟内存管理（paging）
          \item 自主编辑核心记忆
          \item Token 减少 85-93\%
        \end{itemize}

        \vspace{0.2cm}

        \textcolor{darkgreen}{影响力最大}：启发了后续所有框架
      \end{block}

      \vspace{0.2cm}

      \begin{block}{LangChain/LangGraph}
        \footnotesize
        \textbf{记忆模块}：
        \begin{itemize}
          \item 会话缓冲 + 持久化 Store
          \item 命名空间层级组织
          \item 向量检索 + 元数据过滤
        \end{itemize}

        \textcolor{darkgreen}{最广泛采用}
      \end{block}
    \end{column}

    \begin{column}{0.48\textwidth}
      \begin{block}{AutoGen (Microsoft)}
        \footnotesize
        \textbf{多智能体记忆共享}：
        \begin{itemize}
          \item 深度集成 MemGPT
          \item Actor 模式架构
          \item 异步事件驱动
        \end{itemize}

        \textcolor{darkgreen}{企业级首选}
      \end{block}

      \vspace{0.2cm}

      \begin{block}{CrewAI}
        \footnotesize
        \textbf{极简记忆接口}：
        \begin{itemize}
          \item 一键开启：\texttt{memory=True}
          \item 4 种记忆类型自动管理
          \item 支持 Mem0 集成
        \end{itemize}

        \textcolor{darkgreen}{最易上手}
      \end{block}
    \end{column}
  \end{columns}
\end{frame}

\begin{frame}
  \frametitle{新兴趋势：图谱化记忆系统}

  \begin{columns}[T]
    \begin{column}{0.55\textwidth}
      \begin{block}{为什么需要图谱？}
        \footnotesize
        \textbf{向量数据库的局限}：
        \begin{itemize}
          \item 只能找"相似"内容
          \item 丢失实体间关系
          \item 无法多跳推理
          \item 不支持时间演变
        \end{itemize}

        \vspace{0.2cm}

        \textbf{图谱的优势}：
        \begin{itemize}
          \item 保留显式关系链接
          \item 支持路径遍历查询
          \item 追踪知识时间演变
          \item 可解释性强
        \end{itemize}
      \end{block}
    \end{column}

    \begin{column}{0.42\textwidth}
      \begin{block}{代表性工具}
        \footnotesize
        \textbf{Zep/Graphiti}
        \begin{itemize}
          \item 基于 Neo4j
          \item 双时态模型
          \item 94.8\% 准确率
          \item 90\% 延迟降低
        \end{itemize}

        \vspace{0.2cm}

        \textbf{Mem0 Graph}
        \begin{itemize}
          \item 向量 + 图混合检索
          \item 91\% 延迟降低
          \item 90\%+ token 节省
        \end{itemize}
      \end{block}

      \vspace{0.2cm}

      \begin{alertblock}{最佳实践}
        \tiny
        \textbf{混合架构}：向量检索定位候选 + 图遍历返回关系
      \end{alertblock}
    \end{column}
  \end{columns}
\end{frame}

\begin{frame}
  \frametitle{学术前沿：2024-2025 突破}

  \begin{columns}[T]
    \begin{column}{0.48\textwidth}
      \begin{block}{A-MEM (2025)}
        \footnotesize
        \textbf{Zettelkasten 笔记网络}
        \begin{itemize}
          \item 动态链接记忆卡片
          \item 性能翻倍
          \item Token 减少 85-93\%
        \end{itemize}
      \end{block}

      \vspace{0.2cm}

      \begin{block}{HippoRAG (NeurIPS 2024)}
        \footnotesize
        \textbf{海马体索引理论}
        \begin{itemize}
          \item KG 作为联想索引
          \item 多跳 QA 提升 20\%
          \item 成本降低 10-20×
        \end{itemize}
      \end{block}
    \end{column}

    \begin{column}{0.48\textwidth}
      \begin{block}{AgeMem (2026)}
        \footnotesize
        \textbf{强化学习驱动记忆管理}
        \begin{itemize}
          \item 记忆操作作为工具
          \item 三阶段渐进训练
          \item Agent 自主决策记忆
        \end{itemize}
      \end{block}

      \vspace{0.2cm}

      \begin{alertblock}{未来方向}
        \footnotesize
        \begin{itemize}
          \item 多模态记忆（图像、视频）
          \item 多智能体记忆共享
          \item RL 驱动的自主记忆管理
        \end{itemize}
      \end{alertblock}
    \end{column}
  \end{columns}
\end{frame}

\begin{frame}
  \frametitle{业界共识：记忆是 Agent 的核心能力}

  \begin{columns}[T]
    \begin{column}{0.58\textwidth}
      \begin{block}{五种记忆类型（认知科学启发）}
        \footnotesize
        \begin{enumerate}
          \item \textbf{短期/工作记忆}：当前上下文窗口
          \item \textbf{长期记忆}：跨会话持久化存储
          \item \textbf{情景记忆}：具体事件 + 时间戳
          \item \textbf{语义记忆}：事实和知识
          \item \textbf{程序记忆}：技能和规则
        \end{enumerate}
      \end{block}

      \vspace{0.2cm}

      \begin{block}{混合存储成为标准}
        \footnotesize
        \begin{itemize}
          \item 向量数据库：语义相似检索
          \item 知识图谱：关系推理
          \item SQL/NoSQL：结构化数据
          \item 文件系统：原始文档
        \end{itemize}
      \end{block}
    \end{column}

    \begin{column}{0.38\textwidth}
      \begin{alertblock}{关键洞察}
        \footnotesize
        \textbf{记忆架构 > 模型能力}

        \vspace{0.2cm}

        再强大的 LLM，没有记忆系统也：
        \begin{itemize}
          \item 无法积累经验
          \item 无法个性化
          \item 无法持续改进
        \end{itemize}
      \end{alertblock}

      \vspace{0.2cm}

      \begin{block}{对 Coding Agent 启示}
        \tiny
        \begin{itemize}
          \item 必须设计记忆层
          \item 优先考虑图谱存储
          \item 参考 MemGPT 架构
          \item 混合多种存储机制
        \end{itemize}
      \end{block}
    \end{column}
  \end{columns}
\end{frame}

%==============================================================================
\section{实施路线图}
%==============================================================================

\begin{frame}
  \frametitle{分阶段实施路线}

  \begin{block}{阶段 1：MVP（最小可行产品）}
    \footnotesize
    \textbf{目标}：验证 Context Graph 的基本价值

    \begin{itemize}
      \item 实现简单的图结构（Action + State 节点，NEXT 边）
      \item 基于 SWE-bench Lite 测试（300+ 样本）
      \item 对比基线：无记忆 Agent
      \item \textbf{预期}：成功率提升 10-15\%
    \end{itemize}

    \textbf{工作量}：2-3 个月
  \end{block}

  \begin{block}{阶段 2：完整系统}
    \footnotesize
    \textbf{目标}：实现完整的学习与优化机制

    \begin{itemize}
      \item 加入 Insight 节点，实现洞察提取
      \item 实现多种边类型（CAUSES, SIMILAR\_TO, ABSTRACTS）
      \item 在完整 SWE-bench 上测试（2000+ 样本）
      \item 对比 ExpeL、A-Mem 等方法
      \item \textbf{预期}：成功率提升 20-30\%，token 消耗降低 30\%
    \end{itemize}

    \textbf{工作量}：4-6 个月
  \end{block}
\end{frame}

\begin{frame}
  \frametitle{分阶段实施路线（续）}

  \begin{block}{阶段 3：进阶优化}
    \footnotesize
    \textbf{目标}：探索可训练图记忆

    \begin{itemize}
      \item 实现基于 RL 的策略学习
      \item 探索多层图结构（具体→抽象）
      \item 跨项目迁移学习实验
      \item \textbf{预期}：在新项目上零样本/少样本表现优于基线
    \end{itemize}

    \textbf{工作量}：6-9 个月
  \end{block}

  \vspace{0.3cm}

  \begin{block}{阶段 4：扩展应用}
    \footnotesize
    \textbf{目标}：扩展到其他 Coding 任务

    \begin{itemize}
      \item 漏洞检测：结合安全知识图谱（CWE/CVE）
      \item 代码审查：学习项目特定的代码风格和最佳实践
      \item 测试生成：记录有效的测试策略
    \end{itemize}

    \textbf{工作量}：根据具体任务而定
  \end{block}
\end{frame}

\begin{frame}
  \frametitle{技术栈建议}

  \begin{columns}[T]
    \begin{column}{0.48\textwidth}
      \begin{block}{核心技术选型}
        \footnotesize
        \textbf{基础 LLM}
        \begin{itemize}
          \item OpenAI GPT-4
          \item Anthropic Claude
          \item DeepSeek Coder
        \end{itemize}

        \vspace{0.2cm}

        \textbf{图数据库}
        \begin{itemize}
          \item Neo4j（生产级）
          \item NetworkX（原型开发）
        \end{itemize}

        \vspace{0.2cm}

        \textbf{向量数据库}
        \begin{itemize}
          \item FAISS
          \item Qdrant
        \end{itemize}

        \vspace{0.2cm}

        \textbf{代码分析}
        \begin{itemize}
          \item Tree-sitter（AST 解析）
          \item Understand（依赖分析）
        \end{itemize}
      \end{block}
    \end{column}

    \begin{column}{0.48\textwidth}
      \begin{block}{开发框架}
        \footnotesize
        \textbf{Agent 框架}
        \begin{itemize}
          \item LangChain / LangGraph
          \item AutoGPT（参考）
          \item 自研（更灵活）
        \end{itemize}

        \vspace{0.2cm}

        \textbf{实验管理}
        \begin{itemize}
          \item Weights \& Biases
          \item MLflow
        \end{itemize}

        \vspace{0.2cm}

        \textbf{代码执行环境}
        \begin{itemize}
          \item Docker（隔离环境）
          \item SWE-bench 官方工具链
        \end{itemize}
      \end{block}

      \vspace{0.2cm}

      \begin{alertblock}{建议}
        \tiny
        初期用 NetworkX + LangChain 快速验证，成熟后迁移到 Neo4j
      \end{alertblock}
    \end{column}
  \end{columns}
\end{frame}

%==============================================================================
\section{开放问题与挑战}
%==============================================================================

\begin{frame}
  \frametitle{需要解决的核心问题}

  \begin{columns}[T]
    \begin{column}{0.32\textwidth}
      \begin{block}{图设计问题}
        \footnotesize
        \textbf{Q1}：节点粒度如何选择？
        \begin{itemize}
          \item 每个 action？
          \item 还是抽象成高层策略？
        \end{itemize}

        \vspace{0.2cm}

        \textbf{Q2}：如何混合多种边关系？
        \begin{itemize}
          \item 时序 vs 因果 vs 相似
          \item 权重如何设置？
        \end{itemize}

        \vspace{0.2cm}

        \textbf{Q3}：图增长如何控制？
        \begin{itemize}
          \item 如何避免无限膨胀？
          \item 遗忘机制设计
        \end{itemize}
      \end{block}
    \end{column}

    \begin{column}{0.32\textwidth}
      \begin{block}{学习问题}
        \footnotesize
        \textbf{Q4}：检索效率优化
        \begin{itemize}
          \item 大图如何快速检索？
          \item 索引策略
        \end{itemize}

        \vspace{0.2cm}

        \textbf{Q5}：如何定义奖励信号？
        \begin{itemize}
          \item 任务成功/失败
          \item 步数、token 消耗
          \item 如何权衡？
        \end{itemize}

        \vspace{0.2cm}

        \textbf{Q6}：探索-利用平衡
        \begin{itemize}
          \item 过度依赖历史
          \item 错过新方案
        \end{itemize}
      \end{block}
    \end{column}

    \begin{column}{0.32\textwidth}
      \begin{block}{Coding 特定}
        \footnotesize
        \textbf{Q7}：Bug 表示与相似度
        \begin{itemize}
          \item 语义特征提取
          \item 相似度计算
        \end{itemize}

        \vspace{0.2cm}

        \textbf{Q8}：代码版本管理
        \begin{itemize}
          \item 经验失效检测
          \item 版本兼容性
        \end{itemize}

        \vspace{0.2cm}

        \textbf{Q9}：评估指标
        \begin{itemize}
          \item 如何衡量图的价值？
          \item 超越 pass@k
        \end{itemize}
      \end{block}
    \end{column}
  \end{columns}
\end{frame}

\begin{frame}
  \frametitle{风险与应对}

  \begin{block}{风险 1：图构建成本过高}
    \small
    \textbf{问题}：每次任务都要更新图，LLM 调用成本高

    \textbf{应对}：
    \begin{itemize}
      \item 使用小模型（如 GPT-3.5）做初步提取
      \item 批量处理，定期而非实时更新
      \item 只对重要任务提取 insights
    \end{itemize}
  \end{block}

  \begin{columns}[T]
    \begin{column}{0.48\textwidth}
      \begin{block}{风险 2：效果不明显}
        \footnotesize
        \textbf{问题}：图的帮助可能不如预期

        \textbf{应对}：
        \begin{itemize}
          \item 从简单任务开始验证
          \item 分析具体案例，找出失败原因
          \item 对比不同检索策略
        \end{itemize}
      \end{block}
    \end{column}

    \begin{column}{0.48\textwidth}
      \begin{block}{风险 3：过拟合历史经验}
        \footnotesize
        \textbf{问题}：过度依赖历史，僵化

        \textbf{应对}：
        \begin{itemize}
          \item 设置探索概率
          \item 定期清理低置信度经验
          \item 引入随机性
        \end{itemize}
      \end{block}
    \end{column}
  \end{columns}
\end{frame}

%==============================================================================
\section{总结}
%==============================================================================

\begin{frame}
  \frametitle{总结：开发 Coding Agent 的关键要素}

  \begin{columns}[T]
    \begin{column}{0.48\textwidth}
      \begin{block}{必须做的工作}
        \footnotesize
        \begin{enumerate}
          \item \textbf{数据结构}：设计节点和边类型
          \item \textbf{图构建}：轨迹记录 + 洞察提取
          \item \textbf{检索推理}：相似任务 + 路径规划
          \item \textbf{主循环}：决策-执行-更新
          \item \textbf{学习机制}：规则提取 + 统计更新
          \item \textbf{评估体系}：成功率 + 效率 + 学习能力
          \item \textbf{特殊处理}：代码变化 + Bug 相似度
        \end{enumerate}
      \end{block}
    \end{column}

    \begin{column}{0.48\textwidth}
      \begin{block}{差异化优势来源}
        \footnotesize
        \begin{itemize}
          \item \textcolor{darkgreen}{\textbf{记忆层设计}}：图结构 > RAG
          \item \textcolor{darkgreen}{\textbf{失败学习}}：避免重复犯错
          \item \textcolor{darkgreen}{\textbf{跨任务迁移}}：经验复用
          \item \textcolor{darkgreen}{\textbf{可解释性}}：路径可视化
        \end{itemize}
      \end{block}

      \vspace{0.2cm}

      \begin{block}{推荐起点}
        \footnotesize
        \begin{itemize}
          \item \textbf{任务}：Bug 修复 (SWE-bench)
          \item \textbf{MVP}：简单图 + 基本检索
          \item \textbf{对比}：无记忆 baseline
        \end{itemize}
      \end{block}

      \vspace{0.2cm}

      \begin{alertblock}{核心洞察}
        \tiny
        \textbf{记忆是 Agent 的核心能力}——图结构化记忆比扁平 RAG 更适合 Coding 任务
      \end{alertblock}
    \end{column}
  \end{columns}
\end{frame}

\begin{frame}
  \frametitle{参考资源}

  \begin{columns}[T]
    \begin{column}{0.48\textwidth}
      \begin{block}{关键论文}
        \footnotesize
        \textbf{GraphRAG 系列}
        \begin{itemize}
          \item Microsoft GraphRAG
          \item HippoRAG (NeurIPS 2024)
          \item CGR³
        \end{itemize}

        \vspace{0.2cm}

        \textbf{Agent Memory}
        \begin{itemize}
          \item MemGPT (2023)
          \item ExpeL (AAAI 2024)
          \item A-Mem (NeurIPS 2025)
          \item AgeMem (2026)
        \end{itemize}
      \end{block}

      \vspace{0.2cm}

      \begin{block}{开源框架}
        \footnotesize
        \begin{itemize}
          \item Letta (MemGPT 官方)
          \item LangChain/LangGraph
          \item AutoGen (Microsoft)
          \item CrewAI
        \end{itemize}
      \end{block}
    \end{column}

    \begin{column}{0.48\textwidth}
      \begin{block}{记忆工具}
        \footnotesize
        \textbf{图谱化记忆}
        \begin{itemize}
          \item Zep/Graphiti (Neo4j)
          \item Mem0 Graph
          \item TrustGraph
        \end{itemize}

        \vspace{0.2cm}

        \textbf{向量数据库}
        \begin{itemize}
          \item Pinecone
          \item Chroma
          \item Weaviate
        \end{itemize}
      \end{block}

      \vspace{0.2cm}

      \begin{block}{Benchmark}
        \footnotesize
        \begin{itemize}
          \item SWE-bench (Bug 修复)
          \item HumanEval (代码生成)
          \item LongMemEval (记忆评估)
        \end{itemize}
      \end{block}
    \end{column}
  \end{columns}

  \vspace{0.3cm}

  \begin{center}
    \textbf{详细文献列表请参考 report.tex}
  \end{center}
\end{frame}

% biblatex bibliography
\begin{frame}[allowframebreaks]
  \frametitle{参考文献}
  \printbibliography
\end{frame}

\end{document}
